\documentclass[amsmath,amssymb,aps,prd,10pt,twocolumn,showkeys]{revtex4}
\usepackage{graphicx}
\usepackage{mathtools}
\usepackage{verbatim}
\DeclareMathOperator\erfc{erfc}
\DeclareMathOperator\erf{erf}
\DeclareMathOperator{\sgn}{sgn}
\DeclareMathOperator{\snr}{SNR}
\begin{document}

\title{Complex Analysis of Askaryan Radiation: UHECR Reconstruction with Askaryan Radio Array}

\author{Jordan C. Hanson}
\email{jhanson2@whittier.edu}
%\affiliation{Department of Physics and Astronomy, Whittier College}
\author{Damian Iba\~{n}ez-Rodriguez}
\affiliation{Department of Physics and Astronomy, Whittier College}
\date{\today}

\keywords{Ultra-high energy neutrino; Askaryan radiation; Mathematical physics}

\maketitle

\section{Responses for Report of Referee A -- DQ14522/Hanson}

The manuscript presents an analytical model that can describe Askaryan radiation from ultra-high energy cosmic ray (UHECR) cascade in ice and as detected in the Askaryan Radio Array antennas. The work builds upon the previously published references 13, 14 and 15 with additional work, mainly focusing on modeling UHECR radio signal. This work could lead toward successful identification of UHECR, a major source of background in detection of ultra-high energy (UHE) neutrinos using in-ice radio techniques, and hence has significance to the field. I have listed my comments below:

\begin{enumerate}
\item In abstract, ``When UHECR cascades enter solid, RF transparent matter at altitudes where the cascade develops, Askaryan radiation can propagate through the solid matter to RF detectors'' Does it refer only to Askaryan radiation produced by cascade developing in the solid medium or also the Askaryan radiation produced in the air? \\

\textit{Our model refers to Askaryan radiation produced within a medium with constant index of refraction.  In this case, we assume the medium is the firn region of the ice sheet.  This is clarified in the abstract, and in subsequent sections of the manuscript.}

\item In introduction, there is ``...the negative excess charge radiates collectively in the 0.1-1 GHz bandwidth''. In abstract, there is ``radio frequency (RF) pulses in the 0.05-1 GHz bandwidth''. Could author elaborate the difference in the lower bandwidth limit? \\ 

\textit{This is more of a typo than a statement about actual bandwidth.  Technically, the Askaryan spectral amplitude decreases linearly with decreasing frequency, until thermal noise takes over in the 10-50 MHz range.  However, RF noise from the galaxy extends to 80-100 MHz, and detectors like ARA often install high-pass filters with 3 dB frequencies around 100 MHz.  We think 100 MHz, or 0.1 GHz, is the more appropriate number.  This has been changed in a consistent fashion throughout the manuscript.}

\item In introduction, there is ``Due to the unique location of ARA at the South Pole, the contribution from the Askaryan effect dominates the observed RF pulses.'' The explanation is later in Section III (``At 90 degrees S, geomagnetic radiation is expected to be horizontally polarized due to the vertical orientation of the local magnetic field. The vertically polarized (VPol) RF channels are there fore less sensitive to geomagnetic radiation.''). Could author put a hint when it is mentioned first time that the explanation will come up later text or put the explanation closer together? \\ 

\textit{This is a good point.  We have followed this advice by including an explanation in both parts of the text.}

\item In Fig. 7 right plot, I will suggest putting sigma = 0.6 ns line (thick black line) on top of gray lines (or use different color scheme) so it does not get buried in the plot and is easier to spot. What among peak, width, phase or correlation of the lines (UHECR data line and sigma = 0.6 ns model) also being matched in the plot? The peak to peak amplitude (normalized) between UHECR line and sigma=0.6 line has some offset (10 percent). UHECR data oscillates but sigma=0.6 ns model stays flat in either direction of central peak. Could author add comment on that? \\ 

\textit{First, we address the suggestions about the graph format.  We expressed the model lines as a region with $\sigma_t = 0.6\pm0.1$ ns.  That is, the three lines for 0.5, 0.6, and 0.7 ns have been plotted in black, with the label $\sigma_t = 0.6\pm0.1$ ns.  This hopefully reflects more clearly the idea that we are trying to plot a model that has a parameter with an uncertainty.  Second, we address the question about fitting.  The reviewer asked ``What among peak, width, phase or correlation of the lines ... [is] matched in the plot?''  In the first paragraph of Sec. V, we wrote that $\sigma_t$ value is obtained from the fit to the CSW shown in Fig. 4 (bottom right).  In Fig. 7 (right), \textbf{there is no fit}, which is what is remarkable. We have added sentences in Sec. V, first paragraph, emphasizing this. In the first paragraph of Sec. V, we also wrote ``The data was taken from Fig. 2 of [4], ssmoothed with a running average filter, normalized to the same value as the model, and inverted.  Our technique (fitting Eq. 7 to CSWs from raw, observed voltage traces) produces an $\vec{E}$-field field that matches the deconvolution used to produce the reconstructed $\vec{E}$-field in Fig. 2 of [4].''  Third, we address the question about the amplitude and oscillations.  The question of amplitude is addressed by our re-working of the graph.  We only trust our $\sigma_t$ value to a precision of 0.1 ns, and this uncertainty causes our $\vec{E}$-field region to encompass the normalized amplitude.  As for the oscillations, this is a common side effect of deconvolution in the presence of a notch-filter.  In reference [4], the authors write that the \textbf{phase response} was deconvolved.  They write nothing about the amplitude response.  Further, the oscillations in Fig. 7 (right) have a period of a 2-5 nanoseconds, which matches the notch frequency in Fig. 2 of [4], 450 MHz.  This explanation is now included in the second paragraph of Sec. V.}

\item In Fig 4, bottom right (ID 2716-58611) does not have reflected radio pulse while other plots in Fig 4 likely have both primary and reflected pulse. Is there a feature in the plot itself that hints at having just a primary pulse or both primary and reflected pulse? Also, does the delay between primary and reflected pulse of around 20 ns have a physical meaning (eg, what is the source of reflection and how far the reflecting interface is from antenna)? \\

\textit{First, we address the question of Fig. 4 bottom right (2716-58611).  We note that it was not our judgement call to decide that there is no reflection.  Rather, we attempted to fit for a reflection for this event in the same fashion as the others, and the fit never converges.  This is now mentioned in the text.  Second, we address the question of the interpretation of reflections.  Reflections like these in the South polar firn with $n=1.34$ yeids a one-way distance of 2-3 meters.  Slight impedance mismatches between RF connectors, cables, and other RF antennas are common causes of such reflections.  The fact that our $\sigma_{t,2}$ values are larger than our $\sigma_{t,1}$ values suggests a low-frequency reflector.  One candidate is an RF antenna that is not the main antenna in event.  However, this is speculative, so we chose initially not to interpret the reflections.  We only argue that they do not change our interpretation of a UHECR origin for the events.  Reviewer B asked the same question below, but also notes that the reflections do not change the interpretation.  In summary, we have added two sentences in the final paragraph of Sec. IV to address our decision not to interpret the reflections, and to note that they are common in RF systems.}

\item In Figure 6 and other similar plots, it may help to zoom in further (eg between 75 and 175 ns) / use different color scheme, to see the match between model and CR event better. \\ 

\textit{We felt it was important to fix the time-axis for the zoom to 200 ns consistently.  At a sampling rate of 1.5 GHz, this means 25 percent of the plot is about equal to the waveform width.  This does two things: it leaves room for the inset, and it shows the detailed oscillations of the model and data.  The inset provides transparency to skeptical readers.  If we show the whole observed waveform, the reader cannot assume we are hiding something.}

\item In section III, when it says ``... and geomagnetic radiation usually dominates the total electromagnetic pulse when the detector lies far beneath the cascade''.  Could you please point me to which text/plot in reference 3 and 5 is the referred from? \\

\textit{We did not actually have a specific figure in mind when comparing our results to those from the ARIANNA and RNO-G.  Rather, our point is that the ARIANNA observations of the UHECRs were geomagnetically dominated because they were observed at a latitude where the magnetic field is not vertical, and at an altitude far beneath the air shower formation.  The RNO-G collaboration advances the same claim in their publication: the observed events are geomagnetically dominated.  The ARA events, however, are observed at very high latitude and high altitude.  We have re-worked the final paragraph of section III to include more detail on this point.}

\item Data from HPol RF channels is not included in this analysis. Could the conclusion be different if HPol RF is also included in the analysis? How different is the HPol signals from VPol signals for the event geometry considered in this analysis? Could author add some comments about this in the conclusion? \\

\textit{We did not yet have access to the processed HPol data at the time of writing.  The corresponding author is an ARA collaboration member, and the processed VPol data were obtained through a private communication after it was published in ref. [4].  Since the HPol RF channel data were not included ref. [4], we should not speculate on it.}

\end{enumerate}

----------------------------------------------------------------------------------------------------------------

\section{Responses for Report of Referee B -- DQ14522/Hanson}

The authors compare the experimental time-dependent voltage signals in the ARA dipoles at the South Pole generated by cosmic ray air shower cores that strike the snow surface  to the time dependent waveforms predicted by their analytical model described in a previous paper.

This paper looks to be in an important corroboration of the model, but I do suggest several modifications.

\begin{enumerate}



\item The authors should describe how the ARA waveforms in Fig 3-6 were obtained.  In the current draft, at the very end of the paper, the authors describe that the ARA waveform data was directly obtained from the ARA collaboration.  But I think a reference to a private communication the first time the ARA waveform data is introduced will improve the readability.



\item It may help the reader if the time region of the  ``reflected pulse'' is identified in the waveforms shown in one of the figures in figures 3-6. Also, what is responsible for generating the reflected pulse in the ARA signals? I think the authors did a very nice job of showing that the reflected pulses do not affect the results significantly.\\

\textit{Table 2 column 4 'delay between primary and reflected pulses $\Delta t$' can be used on figures 3-6 to identify when a reflected pulse roughly begins. As of now there is no definitive answer to what is causing the reflected pulses in the literature. We could offer some of the speculations; however, what the final answer might be is not impacting our mathematical framework. Though it is agreeable that we have left quite a bit to inference. We append to the final paragraph of section IV our justification.}



\item The paragraph beginning “Several points…” on page 3, argues that pulser calibration events produce a reliable value for f0, but the authors should elaborate on why the calibration events should produce a reliable estimate of gamma, the decay time of the antenna response.  Why doesn’t the double convolution from a pulser signal produce a value for gamma in this analysis that is too large?  Also, the quoted error on the check is larger than the mean, so the authors should argue why this is a meaningful check.



\item The ARA events are different from what is expected from a neutrino interaction in the ice.  The ARA events are generated by the core of a cosmic ray shower that strikes the surface of the snow.  Therefore, the electron distribution is spread over much larger radius (the Moliere radius in air is larger than in ice),  which effects the bandwidth of frequencies that coherently add, and the shower enters the snow surface and encounters a varying index of refraction through the remainder of the shower development.   Both of these effects will alter the time and angular emission dependence  of the Askaryan electric field.  The authors should comment on these differences between the idealized model of neutrino emission assumed in their earlier papers and emission in this case from air shower cores. Given the differences, might the reader expect that the author’s model should provide even better results for neutrino emission within the ice?\\

\textit{Yes, it is a very good inference that our model would work even better with an in-ice neutrino induced Askaryan cascade. An in-ice neutrino induced Askaryan cascade would be cleaner from noise than a UHCR airshower for the reasons listed by you. Our model assumes a consistent index of refraction and does not account for the excess radiation around the air-shower core. A paragraph added to our conclusion section has been added in agreement.}



\item Fig 3:  what is the purpose for a smaller insert graphic?  What additional information is revealed?  Also, the figure caption should explain the insert graph, instead of letting the reader try to figure it out.\\

\textit{The smaller graphic is to display the entire wave form for the sake of transparency to any skeptical critics of the model. A sentence before Fig 3. has been added to create a small break between sections 4 and 5}



\item Fig. 3:  The authors should explicitly state in the figure caption that the model amplitude is relative voltage amplitude from equation 7, and the ARA event amplitude is relative voltage amplitude.  (to avoid possible confusion with the electric field). \\

\textit{Agreed, this is now appended to the aforementioned sentence to clarify the meaning of the inserts and vertical-axis.}



\item In the conclusion, the authors estimate c=0.3/1.78 m/ns, but the index of refraction n=1.78 is value  for ice at depths below ~100m.  The signals from cosmic ray cores are generated near the surface, where index of refraction is much smaller, closer to n=1.4,  and then increases with depth.  Was the estimate in the paper a typo?  If not, the authors should explain the estimate in greater detail.   There is a similar issue with the estimate for the Cherenkov angle.   My last comment is that the parameter a is in units of meters, but since the snow density near the surface is less than half the ice density (0.92 g/cm3) and varies, it will help the reader
to explain if the estimate of a was for ice density, or the snow near the surface?  Usually, lengths in cascade development is parameterized in terms of kg/m$^2$ to account for varying densities of materials.



\item Typo: Measrure on page 7 \\

\textit{Spelling Corrected}

\item Overall assessment of paper.  This paper provides the first experimental proof that Askaryan radio pulses generated by air shower cores developing within the first 10m of the snow surface are described by the analytic model developed in earlier papers, and 
therefore, an important paper to publish. \\



\end{enumerate}

\end{document}
